\chapter{Connection to Perfect Matching}

This chapter describes how a part of a sudoku state can be converted to a bipartite graph and how this can be used to
find strategies: Naked Candidates, Hidden Candidates, and Generalised X-wing with the smallest strategy size $k$.

\section{Preliminaries}

For a bipartite graph $G = (A \cup B, E)$ and $S \subseteq A$,
$N(S)$ denotes the set of neighbours of $S$ in $B$.

A perfectly matchable edge is an edge that can be part of at least one perfect matching.

\begin{thm}[Hall's Marriage Theorem]\label{thm:hall}
    A bipartite graph $G = (A \cup B, E)$ has a perfect matching if and only if
    $\forall S \subseteq A: \left|N(S)\right| \geq \left|S\right|$.
\end{thm}

(TODO: source)

\section{Transforming a Sudoku State to a Bipartite Graph}

This section explains how to transform a part of a sudoku state to a bipartite graph and how
a valid sudoku solution corresponds to a perfect matching of that graph.

\subsection*{Naked candidates}

Let $R$ be a sudoku region and $U$ a state.

We construct a bipartite graph $G := (A \cup B, E)$, where
$A := (A_1, A_2, ..., A_n)$, $B := (B_1, B_2, ..., B_n)$
and $E := \{(A_i, B_j) \mid i, j \in [n] \land j \in U_{R_i}\}$.

A vertex $A_i$ corresponds to $R_i$ --- the $i$-th cell of the region $R$.
A vertex $B_j$ corresponds to the sudoku number $j$.
An edge $(A_i, B_j)$ exists if $j \in U_{R_i}$ --- the number $j$ can be in the $i$-th cell of region $R$.

If the state $U$ is a valid solution, then the graph $G$ is a perfect matching.
(TODO: explain more?, proof?)
Therefore, if an edge $(A_i, B_j) \in E$ is not perfectly matchable,
then there is no valid solution $V \preceq U$ such that $j \in V_{R_i}$.

\subsection*{Hidden candidates}

Let $R$ be a sudoku region and $U$ a state.

We construct a bipartite graph $G := (A \cup B, E)$, where
$A := (A_1, A_2, ..., A_n)$, $B := (B_1, B_2, ..., B_n)$
and $E := \{(A_j, B_i) \mid i, j \in [n] \land j \in U_{R_i}\}$.

Vertex $B_i$ corresponds to $R_i$ -- the $i$-th cell of the region $R$.
Vertex $A_j$ corresponds to the sudoku number $j$.
An edge $(A_j, B_i)$ exists if $j \in U_{R_i}$ -- the number $j$ can be in the $i$-th cell of region $R$.

The construction of the graph $G$ is very similar to Naked candidates strategy.
If $G_N = (A_N \cup B_N, E_N)$ is the graph $G$ for Naked candidates
and $G_H = (A_H \cup B_H, E_H)$ be the graph $G$ for Hidden candidates for the same region $R$ and state $U$,
then $A_N$ and $B_H$ both correspond to the same cells -- cells of $R$,
$B_N$ and $A_H$ both correspond to the sudoku numbers and
$E_H = \{(A_i, B_j) \mid (A_j, B_i) \in E_N\}$.

Similar to $G_N$, if the state $U$ is a valid solution, then $G_H$ is a perfect matching.
(TODO: explain more?, proof?)
Therefore, if an edge $(A_i, B_j) \in E$ is not perfectly matchable,
then there is no valid solution $V \preceq U$ such that $i \in V_{R_j}$.

\subsection*{Generalised X-Wing}

Let $m$ be a sudoku number, $U$ be a state
and $P$ and $Q$ be two sets of $n$ regions such that each set of regions contains every cell of sudoku exactly once

We construct a bipartite graph $G := (A \cup B, E)$, where
$A := (A_1, A_2, ..., A_n)$, $B := (B_1, B_2, ..., B_n)$
and $E := \{(A_i, B_j) \mid i, j \in [n] \land \exists c \in P_i \cap Q_j: m \in U_c\}$.

Vertex $A_i$ corresponds to $P_i$ -- the $i$-th region in $P$.
Vertex $B_j$ corresponds to $Q_j$ -- the $j$-th region in $Q$.
An edge $(A_i, B_j)$ exists if $\exists c \in P_i \cap Q_j: m \in U_c$
-- the number $m$ can be in any cell in the intersection $P_i \cap Q_j$.

Similar to the previous two subsections, if the state $U$ is a valid solution, then $G$ is a perfect matching.
(TODO: explain more?, proof?)
Therefore, if an edge $(A_i, B_j) \in E$ is not perfectly matchable,
then there is no valid solution $V \preceq U$ such that $\exists c \in P_i \cap Q_j: m \in V_c$.

\section{TODO: name}

\begin{thm}[Correctness of Perfect Matching Strategies]\label{thm:strategy}
    Let $G = (A \cup B, E)$ be a bipartite graph, $S \subseteq A$ and $\left|N(S)\right| = \left|S\right|$.
    Any edge $(a, b) \in E$ where $a \notin S$ and $b \in N(S)$ is not perfectly matchable.
\end{thm}

\begin{proof}
%    Let $G' := G - \{a, b\}$ be the bipartite graph obtained by removing vertices $a$ and $b$ from $G$,
%    $A' := A \setminus \{a\}$, $B' := B \setminus \{b\}$
%    and for $S \subseteq A'$, $N'(S)$ denotes the set of neighbours of $S$ in $B'$.
%
%    For a contradiction, let us assume that $(a, b)$ is part of a perfect matching $M$ in $G$.
%    Then $G'$ has a perfect matching $M \setminus (a, b)$.
%
%    Since $\left|N(S)\right| = \left|S\right|$, $a \notin S$, and $b \in N(S)$, then $N'(S) < S$.
%    By Theorem~\ref{thm:hall}, $G'$ does not have a perfect matching.
%    This contradicts $G'$ having a perfect matching $M \setminus (a, b)$,
%    hence $(a, b)$ is not perfectly matchable in $G$.

    All vertices in $S$ can only be matched with vertices in $N(S)$.
    In a perfect matching all vertices in $S$ have to be matched,
    and since $\left|N(S)\right| = \left|S\right|$, there are no vertices in $N(S)$ left.
    Therefore, no vertex outside $S$ can be matched to a vertex in $N(S)$.
\end{proof}

\begin{thm}[Existence of Perfect Matching Strategies]\label{thm:existence}
    Let $G = (A \cup B, E)$ be a bipartite graph with
    a perfect matching and an edge $(a, b)$, that is not perfectly matchable.
    Then there exists a non-empty $S \subsetneq A$ such that $\left|N(S)\right| = \left|S\right|$,
    $a \notin S$, and $b \in N(S)$.
\end{thm}

\begin{proof}
    Let $G' := G - \{a, b\}$ be the bipartite graph obtained by removing vertices $a$ and $b$ from $G$,
    $A' := A \setminus \{a\}$, $B' := B \setminus \{b\}$
    and for $S \subseteq A'$, $N'(S)$ denotes the set of neighbours of $S$ in $B'$.
    If $G'$ had a perfect matching $M$, then $M \cup \{(a, b)\}$ would be a perfect matching of $G$,
    which contradicts the assumption that $(a, b)$ is not perfectly matchable.
    Since $G'$ does not have a perfect matching, by Theorem~\ref{thm:hall},
    $\exists S \subseteq A': \left|N'(S)\right| < \left|S\right|$.
    But since $G$ has a perfect matching, by Theorem~\ref{thm:hall}, $\left|N(S)\right| \geq \left|S\right|$.
    The only vertex that can be in $N(S)$ and not in $N'(S)$ is $b$.
    Hence, $\left|N(S)\right| = \left|S\right|$ and $b \in N(S)$.
\end{proof}

\section{Transforming Unmatchable Edge to Strategy}

Let $G = (A \cup B, E)$ be a bipartite graph with a perfect matching.

By Theorem~\ref{thm:existence}, if there is an edge $(a, b)$, that is not perfectly matchable,
then there is a non-empty set $S \subsetneq A$ such that $a \notin S$ and $b \in N(S)$.
Then by Theorem~\ref{thm:strategy}, if we know $S$, we can prove, that the edge $(a, b)$ is not perfectly matchable.

This section describes how this can be converted to the three sudoku strategies.

(TODO: add explanation for definitions in section: Transforming State to Graph ?)
\xxx{each of the subsections continues on the subsection in the section Transforming State to Graph}

\subsection*{Naked Candidates}

The set $S$ corresponds to a set of cells $X := \{R_i \mid A_i \in S\}$,
the set $N(S)$ corresponds to a set of sudoku numbers $Y := \{i \mid B_i \in N(S)\}$
and the edge $(a, b) = (A_m, B_c)$, corresponds to a sudoku number $m \in Y$,
that is contained in a cell $c \notin X$, which means $m \in U_c$.

Then since $(a, b)$ is not perfectly matchable, $m \notin (S_{ALL})_c$.
That means that if a set of $k$ cells in a region contain exactly $k$ different numbers,
these numbers can not be in any other cell in the region.

\subsection*{Hidden Candidates}

The set $S$ corresponds to a set of sudoku numbers $X := \{i \mid A_i \in S\}$,
the set $N(S)$ corresponds to a set of sudoku cells $Y := \{R_i \mid B_i \in N(S)\}$
and the edge $(a, b) = (A_c, B_m)$, corresponds to a cell $c \in Y$,
that contains the number $m \notin X$, which means $m \in U_c$.

Then since $(a, b)$ is not perfectly matchable, $m \notin (S_{ALL})_c$.
That means that if a set of $k$ numbers in a region is contained in exactly $k$ different cells,
these cells can not have any other number.

\subsection*{Generalised X-Wing}

The set $S$ corresponds to $P_S := \{P_i \mid i \in S\}$,
the set $N(S)$ corresponds to $Q_S := \{Q_i \mid i \in S\}$
and the edge $(a, b) = (A_p, A_q)$ corresponds to the intersection of $P_p \cap Q_q$, where
$P_p \in P \setminus P_S$ and $Q_q \in Q_S$, that contains the number $m$,
which means $\exists c \in P_p \cap Q_q: m \in U_c$.

Then since $(a, b)$ is not perfectly matchable, $\forall c \in P_p \cap Q_q: m \notin (S_{ALL})_c$.
That means that if a sudoku number $m$ can in $k$ regions in $P$ be only in $k$ regions in $Q$,
then $m$ can not be in those $k$ regions in $Q$ outside the $k$ regions in $P$.

\section{Oriented graph H}\label{sec:graph-h}

This section explains how the graph $G$ in (TODO: link to algorithm)
gets transformed to the graph $H$.

Let $G := \{A \cup B, E_G\}$ be a bipartite graph with a perfect matching $P$.

Oriented graph $H := \{V, E_H\}$, $V := \{V_i \mid i \in [n]\}$,
$E_H := \{ (V_i, V_j) \mid i, j \in [n] \land (A_i, B_j) \in E_G \}$
is constructed from $G$ by first orienting the edges $(A_i, B_j)$ to $A_i \to B_j$ and
then contracting all edges that are between the pairs of vertices matched by the perfect matching $P$.

\oneimage{perfect-matching/graph-h}{
    Conversion of $G$ (left) to $H$ (right) (TODO: change image?)
}{0.4\textwidth}

\section{TODO: name}

Let $G := \{A \cup B, E_G\}$ be a bipartite graph with a perfect matching $P$,
$H := \{V, E_H\}$ be the oriented graph as constructed in Section~\ref{sec:graph-h}
and $C$ be the condensation of $H$.

\begin{lemma}\label{lemma:lemma-1}
    Let $S \subseteq A$ and $\left|N(S)\right| = \left|S\right|$, then $N(S) = P(S)$.
\end{lemma}

\begin{proof}
    Only neighbours of $s \in A$ can be perfectly matched, therefore $N(S) \supseteq P(S)$.

    For a contradiction, suppose that $N(S) \neq P(S)$.
    Since $N(S) \supseteq P(S)$ and $N(S) \neq P(S)$, we have $N(S) \supsetneq P(S)$,
    hence $\left|N(S)\right| > \left|P(S)\right|$.
    And since $\left|P(S)\right| = \left|S\right|$, we have $\left|N(S)\right| > \left|S\right|$,
    which contradicts the assumption that $\left|N(S)\right| = \left|S\right|$.
\end{proof}

\begin{lemma}\label{lemma:lemma-2}
    Let $S \subseteq A$ and $N(S) = P(S)$.
    Suppose that $A_i \in S$, then for all vertices $V_j \in V$, that are reachable from $V_i$, $A_j \in S$.
\end{lemma}

\begin{proof}
    For a contradiction, suppose that there is a vertex $A_i \in S$ and vertex $A_j \notin S$,
    such that $V_j$ is reachable from $V_i$.
    This means that there is a path from $V_i$ to $V_j$ and on that path there are two neighbouring vertices
    $V_x$ and $V_y$ such that $A_x \in S$, $A_y \notin S$ and $(V_x, V_y) \in E_H$.
    Since $(V_x, V_y) \in E_H$, $(A_x, P(A_y)) \in E_G$.
    Hence $A_y \notin S$ and $P(A_y) \in N(S)$ which contradicts $N(S) = P(S)$.
\end{proof}

\begin{lemma}\label{lemma:lemma-3}
    Let $S \subseteq A$ be a non-empty set and $N(S) = P(S)$.
    Suppose that there exists an edge $(a, b) \in E_G$ such that $a \notin S$ and $b \in N(S)$.
    Then there exists a sink $X$, that is not a source, in the condensation $C$ containing vertices $X \subseteq V$,
    such that $\{A_i \mid V_i \in X\} \subseteq S$.
\end{lemma}

\begin{proof}
    Pick a vertex $A_x \in S$ such that $P(A_x) = b$.
    Let $X_1$ be the SCC containing $V_x$.
    By Lemma~\ref{lemma:lemma-2}, if $A_x \in S$, then $\{A_i \mid V_i \in X_1\} \subseteq S$.
    Since $(a, b) \in E_G$, $a \notin S$, and $b \in N(S)$, then $X_1$ is not a source.

    Since $C$ is DAG, $X_1$ can reach at least one sink $X_l$.
    Since $X_1$ is not a source, the sink $X_l$ is not a source either.
    The path $(X_1, X_2, ..., X_l)$ in the condensation $C$ can be converted to a path
    $p := (V_{D_1}, V_{C_2}, ..., V_{D_2}, V_{C_3}, ..., V_{C_l})$,
    where $V_{D_1} = V_x$, $V_{C_l}$ is any vertex in the SCC $X_l$ and
    the pair $V_{D_i}$ and $V_{C_{i+1}}$ is a pair of vertices in SCCs $X_i$ and $X_{i+1}$,
    such that $(V_{D_i}, V_{C_{i+1}}) \in E_H$.
    Since $V_{C_i}$ and $V_{D_i}$ are in the same SCC $X_i$, then there exists a path from $V_{C_i}$ to $V_{D_i}$,
    therefore $p$ is a path in $H$.
    (TODO: is it possible to not have to do all this to prove that a vertex in a sink $X_l$ is reachable from a vertex from $X_1$)

    Since $p$ is a path from $V_x$ to $V_{C_l} \in X_l$, then all vertices $V_i \in X_l$ are reachable from $V_x$,
    which means that $\{A_i \mid V_i \in X_l\} \subseteq S$.
\end{proof}

\begin{lemma}\label{lemma:lemma-4}
    Let $S \subseteq A$ be the smallest non-empty set such that $\left|N(S)\right| = \left|S\right|$ and
    there exists an edge $(a, b) \in E_G$ such that $a \notin S$ and $b \in N(S)$.
    Then $X := \{V_i \mid A_i \in S\}$ is a sink and not a source in $C$.
\end{lemma}

\begin{proof}
    By Lemma~\ref{lemma:lemma-3}, there exists a sink $Y$, that is not a source, in $C$,
    such that $\{A_i \mid V_i \in Y\} \subseteq S$.

    For a contradiction, suppose that $\{A_i \mid V_i \in Y\} \neq S$.
    Then let $S' := \{A_i \mid V_i \in Y\}$.
    Since $\{A_i \mid V_i \in Y\} \subsetneq S$, then $S' \subsetneq S$.
    Since $Y$ is a source, then $S' = P(S')$, hence $\left|N(S')\right| = \left|S\right|$.
    Since $Y$ is not a sink, then there exists an edge $(a, b) \in E_G$ such that $a \notin S'$ and $b \in N(S')$.
    This contradicts the assumption that $S$ is the smallest set such that $\left|N(S)\right| = \left|S\right|$ and
    there exists an edge $(a, b) \in E_G$ such that $a \notin S$ and $b \in N(S)$.
    That means that $\{A_i \mid V_i \in Y\} = S$, then $X = Y$.
\end{proof}

\section{Algorithm}

This section explains the algorithm used to find the smallest set $S$ from Theorem~\ref{thm:strategy}.

\begin{algorithm}
    \caption{Smallest perfect matching strategy}
    \begin{algorithmic}[1]
    \Require Bipartite graph $G = (A \cup B, E)$
    \Ensure If $G$ does not have a perfect matching: \texttt{no-matching}
    \Ensure Else if all edges in $G$ are perfectly matchable: \texttt{no-edge}
    \Ensure Else: The smallest possible $S \subseteq A$ such that $\left|N(S)\right| = \left|S\right|$
    and there exists an edge $(a, b)$, that is not perfectly matchable, $a \notin S$, and $b \in N(S)$.
    \State $P \gets$ a maximum matching of $G$ using Hopcroft-Karp algorithm
    \If{$P$ is not a perfect matching}
        \Return \texttt{no-matching}
    \EndIf
    \State $H \gets$ graph $H$ made from graph $G$ by contracting all perfectly matched edges by $P$
    \State $C \gets$ components of strong connectivity of $H$
    \State $D \gets$ condensation of $C$
    \State $S \gets$ list of sources of $D$ that are not sinks
    \If{$S$ is empty}
        \Return \texttt{no-edge}
    \EndIf
    \Return Smallest element of $S$
    \end{algorithmic}\label{alg:algorithm}
\end{algorithm}

\begin{thm}{Correctness of Algorithm~\ref{alg:algorithm}}
    If the bipartite graph $G$ has a perfect matching and at least one edge $(a, b)$, that is not perfectly matchable,
    then Algorithm~\ref{alg:algorithm} has a correct output.
\end{thm}

\begin{proof}
    Let $G := (A \cup B, E)$ be a bipartite graph with a perfect matching $P$ and at least one edge $(a, b)$,
    that is not perfectly matchable.

    By Theorem~\ref{thm:existence}, there exists a non-empty $S \subsetneq A$,
    such that $\left|N(S)\right| = \left|S\right|$, $a \notin S$, and $b \in N(S)$.
    Then by Theorem~\ref{lemma:lemma-4}, for the smallest non-empty $S$,
    There exists a sink $X := \{V_i \mid A_i \in S\}$, that is not a source, in $C$.

    By searching through all sinks, that are not sources, in $C$, the algorithm finds the smallest set $S$
    and returns it.
\end{proof}
